\section{Introduction}\label{sec:intro}

The Federal Reserve departed from its ``Treasuries-only'' framework in 2008 and 2020, purchasing trillions in agency MBS that suppressed 30-year mortgage rates to lows of roughly 2.6\% \citep{hancock2011} and triggered a refinance boom. In 2022 it reversed course. Quantitative Tightening (QT) was a passive roll-off intended to drain up to \$35 billion per month from its SOMA portfolio, under a cap phased in at \$17.5 billion and stepped up three months later. (June 2024 and April 2025 runoff slowings cut only the Treasury cap, and above-cap reinvestment never triggered in this window. So the constant post-phase-in \$35 billion MBS schedule is policy-committed throughout.) Unlike standard Treasuries, U.S. mortgages carry asymmetric, state-dependent payoff profiles. Tightening pushed the 30-year rate past 7.0\%, and under the par-payoff rule that imposed a substantial penalty on relocation: the lock-in effect. In practice, historic housing-inventory shortages and depressed consumer sentiment compounded it. The rate rise decomposes into the 10-year Treasury yield's climb plus primary-mortgage spread widening from roughly 170 to over 300~bp at its 2022--2023 peak (FRED series MORTGAGE30US less DGS10). Spread components follow \citet{boyarchenko2019}, whose framework predates and does not document the 2022--2023 magnitudes. Mortgages did not prepay. Cumulative tightening-cycle runoff fell \$764.7 billion short of the phased caps \citep[on the balance-sheet setting, see][]{lopezsalido2023,acharya2023} (Section~\ref{sec:method-benchmark}). Sections~\ref{sec:abm} and~\ref{sec:hazard} decompose that shortfall, isolating the lock-in channel's own share.

That the lock-in effect slowed the runoff of the Federal Reserve's mortgage portfolio is not a new claim; its own staff and officials have said as much (\citealp{na2024}; \citealp{perli2024}; \citealp{hammack2025}). This paper's contribution is not that mechanism but its decomposition: an accounting exercise within a fully disclosed model, not a causal identification. Switching the lock-in elasticity off partitions the realized shortfall between an elastic channel and a rate-inelastic baseline: scheduled amortization plus baseline turnover held at its observed level. Its floor is read off realized, partly behavioral turnover, so the baseline is not a non-behavioral object. The elasticity's causal content is inherited from the external estimates calibrating it, not established here (Section~\ref{sec:identification}). Most of the shortfall sits in that baseline: with the elasticity off, the loan-level model still recovers 85.7\% of the \$764.7 billion shortfall-against-cap benchmark (Section~\ref{sec:method-benchmark}) at the calibration this paper headlines and 88.7\% in-sample, on the benchmark-consistent accounting basis (Section~\ref{sec:hazard-interp}). The Federal Reserve's own ex-ante projection implies the same composition, though not independently: it may already have embedded part of the lock-in channel. The New York Fed's May 2022 staff baseline already anticipated the large majority of the realized cap-shortfall: 75.6--88.5\%, depending on a disclosed intra-2022 allocation choice. That range is wide enough that its overlap with the rate-inelastic null's recovery carries no precision.

The elasticity's identified footprint is therefore a bounded margin, not a pinned magnitude and not a sign. The floor that dominates the level is calibrated in-window, on the same locked-in cohorts the margin decomposes. I therefore re-measure it outside that window, on pre-episode 2017--2019 discount cohorts, and re-run both legs at the result. The margin runs $+4.3$ to $+6.8$ points across the off-window floor's defensible range. It runs $+2.3$ to $+9.1$ points once that floor's own sampling error is propagated with a wild-cluster small-sample correction and the pre-committed coverage rule selects the qualifying construction. That is the binding layer (Section~\ref{sec:robustness-floor}). At the mid-grid anchor inside that interval the value is $+5.6$ points, or $+\$42.6$ billion. The design delivers the range, not that value. Measured corrections to the floor and the accounting basis run in both directions: calendar-standardizing the floor read moves the value to $+4.7$ points; matching that read's housing-activity level to the window's moves it to $+10.7$. Specification variants that are not floor corrections run both ways as well, so the identified content sits in the range and not in any interior member (Section~\ref{sec:identification}). The in-window production calibration returns $+9.2$ points, reported as the in-sample calibration point, not the headline. Its sampling intervals are in Table~\ref{tab:headline}, and the operative uncertainty remains calibration (Table~\ref{tab:uncertainty}). Neither figure is the entire shortfall (Section~\ref{sec:identification}).

A pre-committed sweep of the baseline turnover floor and the elasticity band bounds the margin between $+2.1$ and $+13.2$ points. The margin stays positive at every point of the calibration box and at every off-window anchor tested. That positivity checks the specification; it does not identify the mechanism (Section~\ref{sec:identification}). A functional-form test replacing the floor's hard maximum with a competing-risks combination leaves the margin inside that box and nearly floor-invariant (Section~\ref{sec:robustness-floor}). The grid, parity gates, and interpretive thresholds were committed to the repository before any run executed. I reserve ``pre-registered'' for third-party registries. That pre-commitment is a disclosure device, not a proof. The sweep's credential should be read net of the replication package's verdict appendix, which tabulates every such rule together with how its outcome was adjudicated after the run \citep{olken2015}.

I quantify this shortfall's mechanics with two structurally distinct estimators of the aggregate level. The first is a 10,000-household Agent-Based Model with dynamic macroeconomic frictions. The second is a loan-level hazard framework that abandons the household utility function altogether, introduced after the ABM's residual proved persistent and unresponsive to added behavioral realism. They disagree sharply about how much household choice explains (Sections~\ref{sec:abm} and~\ref{sec:hazard}). The ABM recovers 13.6\% on the fifty-seed mean (10--17\% across seed noise; 11.9\% on the frozen production draw), the hazard framework 91.3\% at the headline off-window floor and 97.9\% at the in-sample calibration point (remaining variants: Table~\ref{tab:headline}). The decomposition rests on the hazard framework alone: the rate-inelastic/rate-elastic partition is a hazard-framework object, because Path A is excluded from every headline range by pre-committed specification, and the ABM's null admits two readings that disagree in sign. Path A lands at 112.4\% yet does not corroborate: no useful propagated precision, no inferential object recomputed under its production specification (Section~\ref{sec:patha}). Two qualifiers travel with the two-estimator contrast wherever I state it. The $\beta_1 = 0$ null already recovers 85.7\% (88.7\% in-sample), so the hazard framework's margin over a rate-inelastic model is $+5.6$ points at the off-window floor ($+9.2$ in-sample), not the level (Section~\ref{sec:identification}). In flow terms the margin is roughly \$1 billion per month against a \$35 billion monthly redemption cap. A cross-design variant of the ABM recalibrated on real loan covariates recovers 59.3\% on the frozen seed, 60.2\% across fifty seeds, and 76.3\% at the book's composition (Section~\ref{sec:robustness-crossdesign}). So the contrast is not clean evidence about modeling paradigm alone. The hazard framework is the instrument; the decomposition it yields is the paper's central empirical result. A counterfactual on the Danish market-value repurchase regime further shows the institutional cash-flow component of systemic duration risk to be a consequence of U.S. institutional design. Under the rule-only transplant it matches the lock-in marginal in origin, and in size at the in-sample calibration point on which both were run.

This result sits inside a two-sided exchange. The pool homogeneity behind To-Be-Announced (TBA) market liquidity \citep[the market's institutional design is reviewed in][]{fuster2023} is purchased with household mobility, because the standardized contracts TBA delivery presupposes carry the ``par-payoff'' obligation. Under that obligation, a homeowner must extinguish the existing mortgage at full face value to relocate. The obligation is a property of the U.S. fixed-rate note, not of TBA eligibility. It is enforced by the due-on-sale clause (made broadly enforceable by the Garn--St.~Germain Act of 1982). The FHA/VA loans backing Ginnie Mae pools are assumable, a statutory carve-out whose exercise is underwritten (HUD Handbook 4000.1; 38 U.S.C.~\S~3714). Those pools are 20.4\% of SOMA book face on the paper's CUSIP-level cohort parse. This design nowhere measures their realized take-up, so that 20.4\% share bounds the carve-out from above. The observed speeds bound it from above a second way, in the same direction. An assumed loan does not prepay, so materially high take-up would show up as \emph{slower} Ginnie voluntary speeds. Over the window Ginnie's voluntary component instead runs $1.35$ points per year \emph{above} the conventional one. Section~\ref{sec:pathb} uses that differential to bracket what a Ginnie share that is not gap-inert would contribute, on a mapping it assumes rather than measures. Both bounds are ceilings, not sizes. Neither prices the frictions (second-lien gaps, equity buy-outs, qualification) between the statutory right and its exercise. The decomposition disciplines this frame. Household mobility is constrained, as the household-level lock-in literature documents, while the cash-flow shortfall is governed mostly by loan-level and institutional mechanics. The trade-off's binding cost is therefore denominated in mobility, not institutional cash flow (Section~\ref{sec:conclusion}). Section~\ref{sec:abm-danish}'s Danish counterfactual prices the payoff rule itself: recalibrated against separately estimated Danish elasticities, the rule-only gap is $+\$61.2$ billion (8\% of the benchmark). That figure is the zero-refinance edge of a band swept to $+\$256.8$ billion at 3\% refinance-in-place. The sign is forced rather than found: pinned at the U.S. zero-gap hazard, the Danish leg must outrun a book out of the money almost everywhere. The informative content is the magnitude. The bracketing anchor that also imports Danish baseline mobility reverses it to $-\$99.9$ billion. The cash-haircut incidence reading reverses it to $-\$51.0$ billion at the re-derived discount. The relief is established under face accounting only, the benchmark's own denomination (Section~\ref{sec:abm-danish}; Appendix~\ref{app:params}). The mobility-versus-cash-flow ranking rests on external evidence, not a within-paper comparison. I measure the institutional leg and import the mobility leg at an elasticity I never re-measure, so the two are nowhere commensurable. The one household-side magnitude I carry is \citepos{batzer2024} roughly \$2.4 trillion of below-market financing value forfeited on relocation. It is computed on universal relocation and therefore an upper bound. A stock of retained financing value is also not a cash flow over forty-two months. The ranking's direction is what that supports; its ratio is not.

This paper's affirmative contributions are three. First, a decomposition of the QT mortgage-runoff shortfall into rate-inelastic and rate-elastic components, the elastic part delivered as a bounded margin (Sections~\ref{sec:abm}, \ref{sec:hazard} and~\ref{sec:identification}). Second, an expectations-based complement that measures the shortfall against the Federal Reserve's own ex-ante projection, not the never-binding caps (Section~\ref{sec:method-benchmark}). Third, the cap-design arithmetic. New York Fed staff projected from the month the \$35 billion ceiling took effect that the caps would sit above achievable runoff \citep{nyfed2022}. The observation is theirs; I add its quantification and ex-ante construction. The ceiling was set roughly 1.7 to 1.9 times the Federal Reserve's own contemporaneous projection of achievable runoff. Both components of the rate-inelastic path are computable in advance of the cycle, as a band (Sections~\ref{sec:identification} and~\ref{sec:discussion-capdesign}). Two further design elements are not contributions. The paired cross-design test is a validation device under interpretive thresholds fixed ex ante, and it ran against this paper's own ABM reading: on all fifty seeds the real-covariate leg crossed the pre-fixed threshold at which that reading is undercut (Section~\ref{sec:robustness-crossdesign}). The order-of-magnitude reduction of the Danish market-value counterfactual, with the incidence bracket that prices it, is a robustness result on a counterfactual. It is not a measurement of the U.S. book. The \$925.5 billion figure it corrects is this paper's own earlier extrapolation, and the corrected gap's sign is anchor- and incidence-determined (Section~\ref{sec:abm-danish}; Appendix~\ref{app:params}).

Table~\ref{tab:headline} collects the core results and their operative uncertainties in one place; every value repeats a committed, gated figure.

{\footnotesize
\begin{longtable}{@{}p{4.6cm}p{4.6cm}p{5.3cm}@{}}
\caption{Core results and operative uncertainties. Each row repeats committed values derived in the referenced exhibit; recovery percentages are quoted on the benchmark-consistent shared basis. The uncertainty hierarchy is established in Table~\ref{tab:uncertainty}: simulation-seed noise is negligible at reporting precision and sampling is quoted under the stratum-cluster scheme; the calibration box is the operative uncertainty --- except the headline marginal, whose binding layer is the floor read's own propagated sampling error (Table~\ref{tab:uncertainty}). The replication package ledger's crosswalk table maps each headline quantity to its accounting basis and source run. The Path B recovery is estimated on the Freddie 2017--2021 conventional universe; Section~\ref{sec:method-benchmark} derives that universe's coverage of SOMA book face on the book's own agency $\times$ vintage joint cells, and states the external-speed overlay that bounds each exclusion (Section~\ref{sec:pathb}).}
\label{tab:headline}\\
\toprule
Object & Headline value & Uncertainty / bounds \\
\midrule
\endfirsthead
\toprule
Object & Headline value & Uncertainty / bounds \\
\midrule
\endhead
Cap-relative benchmark, June 2022--November 2025 (Section~\ref{sec:method-benchmark}) & \$764.7 billion & -- \\
Expectations-based complement (Section~\ref{sec:method-benchmark}) & \$87.8 billion, 11.5\% of the cap-relative benchmark & Fed's ex-ante projection implied 75.6--88.5\% of the cap-shortfall, by intra-2022 allocation \\
Rate-inelastic baseline, $\beta_1 = 0$ (Section~\ref{sec:identification}; Table~\ref{tab:bases}) & 85.7\% of benchmark at the headline off-window floor; 88.7\% in-sample, both under the production floor form (35.6\% under the additive form) & 83.9--86.9\% across the off-window floor range under that form; moves one-for-one with the floor level; floor-read wild-cluster $[80.6, 88.4]$\% (run \texttt{null\_floor\_interval}) \\
Lock-in marginal, off-window floor (Sections~\ref{sec:identification} and~\ref{sec:robustness-floor}; Table~\ref{tab:oosfloor}) & $+2.3$ to $+9.1$ points, basis-invariant --- the floor read's own sampling error propagated at the central elasticity (the 6.5\% mid-band value), not an interval on the elasticity, whose dimension Table~\ref{tab:lowband} sweeps; $+5.6$ points $= +\$42.6$ billion at the mid-grid anchor inside it, an anchor convention rather than a central tendency; the range carries the identified content (Section~\ref{sec:identification}) & floor-read wild-cluster bootstrap-$t$ inversion $[+2.3, +9.1]$, the binding layer (31 clusters; percentile read $[+3.0, +8.0]$, under-covering; the corrected ladder and its degrees of freedom are Table~\ref{tab:ladder})\tnote{a} \\
Lock-in marginal, convention envelope (Table~\ref{tab:assembly}) & $+0.9$ to $+15.6$ points across named convention variants (seasoning ramp 75--150 PSA; moving-share bracket; housing-activity-scaled baseline; floor form) & convention ranges without coverage properties; the calibration box spans $+2.1$ to $+13.2$ \\
\addlinespace
Lock-in marginal, in-sample calibration point (Section~\ref{sec:identification}; Table~\ref{tab:uncertainty}) & $+9.2$ points $= +\$70.3$ billion, basis-invariant & sampling: within-stratum $[+9.17, +9.23]$, stratum-cluster $[+8.27, +10.19]$ (the operative scheme); calibration box $+2.1$ to $+13.2$; Fannie replication $+8.68$ (Section~\ref{sec:pathb-fannie}) \\
Path B central level (Table~\ref{tab:estimators}) & 91.3\% at the headline off-window floor; 97.9\% in-sample (100.0\% composed as committed, 98.1\% on a single coupon convention) & 88.2--93.7\% across the off-window floor range; covariate priors 92.2--99.2\% (in-sample); Ginnie-speed overlay at the headline floor 84.0\% central against a 79.5\% null, and 89.3\% at the in-sample point (raw; of the \$66.3B raw adjustment the Ginnie-specific piece is $+\$38.5$B, within the \$20--47B bound) \\
Production ABM (Table~\ref{tab:estimators}) & 13.6\% seed mean & 10--17\% seed band; frozen draw 11.9\%; the operative uncertainty is calibration, not seed noise --- established on the real-covariate cross-design, where one calibration choice moves recovery across 20.9--59.3\% (Section~\ref{sec:robustness-crossdesign}); that span belongs to two calibration choices on that design, not to this row's 13.6\%, and does not contain it \\
Cross-design ABM, real covariates (Table~\ref{tab:crossdesign}) & 59.3\% recalibrated; fifty-seed mean 60.2\%, 95\% band 55.1--66.0\% (76.3\% at book composition; effective sample 1{,}332 of 10{,}000, re-draw 70.6\%) & 20.9\% under the frozen scale; fifty-seed mean 22.0\% (12.6\% at book composition) \\
Danish rule-only counterfactual (Table~\ref{tab:danish}) & $+\$61.2$ to $+\$256.8$ billion over 0--3\% refinance-in-place, of which $+\$61.2$ billion is the zero-refinance edge rather than a central value & positivity is forced by the zero-gap anchor, so the refinancing sweep cannot flip it; the Danish-level bracketing anchor does, at $-\$99.9$ billion. The gap is incidence-conditional: $+\$61.2$ billion under face accounting (the benchmark's own denomination), reversing to $-\$51.0$ billion under cash accounting at a re-derived discount, and product-mix-conditional besides: under face accounting itself it reverses above a Danish interest-only share of 34.3\% (runs \texttt{buyback\_credit\_bracket} and \texttt{buyback\_discount\_rederived}; Appendix~\ref{app:params}) \\
\addlinespace
Household-side foregone payoffs implied by the marginal (Section~\ref{sec:discussion-policy}) & about 179{,}500 on the book over the window (51{,}300 a year); 233{,}000--256{,}000 a year grossed to the whole market & an \emph{implied} count, not a measured one --- the marginal divided by a representative surviving balance, inheriting every conditionality the marginal carries. Upper bound at the production $\mu = 1$ convention (108{,}300 at $\mu = 0.5$, 58{,}700 at $\mu = 0.25$, on the book; the grossed upper bound is 5.7--6.3\% of the 4.08 million annual existing-home-sale run rate; Section~\ref{sec:discussion-policy}). Scaled by the household-level literature's per-move welfare estimates, the counts are the marginal's household-side image, the magnitude attributed to that literature \citep{gerardi2024} rather than measured here. \citet{batzer2024}'s \$2.4 trillion of below-market financing value prices the same households' stakes on universal relocation and is not additive with this count. \citet{hsieh2019} give the misallocation mechanism the moves would feed, not estimated here \\
\addlinespace
Portfolio duration extension (Table~\ref{tab:wal}; Appendix~\ref{sec:robustness-wal}) & empirical-path WAL 9.4 years at the window's open and 8.5 at its close --- a 6.0-year extension against the no-shock counterfactual's 2021 refinancing-boom speeds (3.4 years) & the comparator sets the result, and this row is quoted against both. Read instead against the baseline turnover floor's off-window read, the same path runs 9.4 years against 9.5 at the open and 8.5 against 8.6 at the close, so the six years measure refinancing-versus-turnover rather than lock-in-versus-normalcy. The Danish rule-only leg shortens the portfolio by 0.6 years against the U.S. leg's 9.7 (Danish-level bracketing leg 10.8). Approximate single-pool-per-cell calculation; option-adjusted duration and convexity are declined as out of scope \\
\bottomrule
\end{longtable}
}

\noindent{\footnotesize \emph{Notes to Table~\ref{tab:headline}.} a Sensitivity catalogue for the headline marginal. Each entry is a committed figure derived where it is cited. None is a sampling interval. Off-window floor range $+4.3$ to $+6.8$ (max form; $+3.9$ to $+7.5$ joint with the band); additive form $+11.2$, floor-invariant; form-conditional hull $+3.5$ to $+13.1$, upper end from a 61.2\%-recovery cell and not read as an equally credentialed member. The concave transform and the additive form do not compose additively: jointly $+9.2$, interaction $-1.47$. Mixture curve in the floor's strictly-involuntary share (Section~\ref{sec:robustness-floor}): $+9.7$ at $\omega{=}0.25$, $+10.7$ at $\omega{=}0.4$; Ginnie overlay $+4.4$ (conventional-share scaling $0.797\times$); floor-stability check $+\$45.1$ billion.\par}

\paragraph{Definitions used throughout.} Four accounting bases attach to every recovery percentage (Section~\ref{sec:method-benchmark}): \emph{standalone-scorer} (no adjustment), \emph{benchmark-consistent shared} (nets a common \$69.6 billion curtailment flow identically from every U.S. leg; the main-text convention), \emph{full-book} (reweights to the SOMA book's coupon mix), and \emph{composed} (both). The realized SOMA CPR is a back-out and moves with the coupon convention it is backed out at: unless a site says otherwise, every empirical CPR here is the \emph{ABM-basis} back-out---window mean 5.14\%, against a hazard-basis 5.52\% and 5.79\% at the note-rate WAC, the corrected amortization basis since the note rate is what amortizes the loan (run \texttt{coupon\_convention\_amortization}; duration pricing in Table~\ref{tab:wal}). Two floor calibrations attach to every marginal and level: the \emph{in-sample calibration point} (the 4\% baseline turnover floor, in-window on the locked-in 2023--24 cohorts; returns $+9.2$ points) and the \emph{off-window calibration} (the 4.70--5.33\% floor, on pre-episode 2018 discount cohorts; returns $+5.6$ points and is the headline). The floor is a measured object: realized-turnover reads put it at 3.8--3.9\% in-window (flat across seasoning cuts); 4.70--5.33\% on the off-window 2018 leg, all of it on cohort-months aged twelve to twenty-four, so that leg's mature-seasoning test is not computable rather than passed; 5.51\% under an imputation-heavy age standardization (84\% imputed weight, an indicative bound from the same 2018 leg); and 5.52\% on an out-of-agency Fannie Mae read of the pooled off-window period---the last two sit above the clean band, as does every mature-versus-young turnover contrast the data make measurable, hence the band's soft lower edge (Section~\ref{sec:robustness-floor}). The ``baseline turnover floor'' (sometimes called an involuntary-turnover floor) names that object without decomposing it: the reads measure total deep-discount-cohort turnover---discretionary life-cycle moves and cash-out refinancings mixed with strictly involuntary events---and the strictly-involuntary share, plausibly well under half (Section~\ref{sec:pathb}), is priced directly by the mixture curve of Section~\ref{sec:robustness-floor}. ``Trapped liquidity'' and ``shortfall'' denote the net shortfall of realized roll-off against the phased redemption caps (\$764.7 billion, June 2022--November 2025). The \emph{lock-in marginal}, the paper's identified object, is the central simulation minus the \emph{rate-inelastic baseline}, its $\beta_1 = 0$ null: scheduled amortization from the book's own coupon, age, and term composition plus baseline turnover held at its observed level, the marginal being the imported elasticity's increment above it. The partition separates baseline from increment, not mechanics from behavior---the floor is read off realized, partly behavioral turnover, so the baseline is not a non-behavioral object and I do not call it one. The marginal's operative uncertainty is the calibration-and-form envelope of Table~\ref{tab:uncertainty} rather than any sampling interval, and the envelope is a grid, not a confidence region---a load-bearing distinction: the elasticity band it sweeps is \citeauthor{liebersohn2024}'s \emph{specification} range with my adopted midpoint at its centre, not a published point estimate with a published standard error, no sampling error of theirs is propagated anywhere in this paper, and so the envelope has no coverage property. Where a genuine sampling interval exists---the floor reads' cluster bootstrap, the loan- and stratum-level bootstraps of Section~\ref{sec:pathb}---I quote it as such and rank it against the envelope explicitly rather than merging the two. One global property is stated here as well as in the limitations, because every number below inherits it: \emph{no outcome-holdout months exist anywhere in this paper}, without exception; its unavoidability is scoped to the cash-flow margin. The identified object has nothing to hold out---a difference between two simulations answers to no realized outcome---while a level-fit holdout, conceivable but not attempted, would split the benchmark itself, one cumulative path against one cap schedule; the household margin does admit outcome moments this design leaves unused, importing the mobility elasticity rather than estimating it (Section~\ref{sec:limitations}). The floor that sets Path~B's level is calibrated on in-window turnover, Path~A trains on nineteen of the forty-two QT months, and Path~B's level is evaluated on the window that sets it---each disclosed again where it arises, and the nineteen-month floor variant reported later is an input-stability check rather than an outcome holdout: what it holds out is calibration months, not the benchmark. Every recovery percentage is accordingly in-sample or calibrated in the time dimension---a description of fit rather than a prediction (Section~\ref{sec:limitations}).
